%% ============================================================
%%  CHAPITRE 3 — Stratégie d'orchestration et observabilité
%%               de la couche d'agents
%% ============================================================
\chapter{Stratégie d'orchestration et observabilité de la couche d'agents}
\label{chap:orchestration}

%% Listes d'outils \mcptool{...} denses : autoriser les coupures à _ , ( etc.
\sloppy

Le chapitre~2 s'est achevé sur un constat à trois faces. La
\emph{propriété de l'orchestration} — selon qu'un superviseur \gls{llm}
ou du code déterministe décide du prochain pas — est l'axe qui
gouverne le plus directement l'auditabilité et la reproductibilité
d'une variante. Les variantes diffèrent ensuite moins par la qualité
intrinsèque du contenu produit que par la \emph{posture humaine}
qu'elles supposent : enseignant disponible en temps réel, opérateur
qui pilote une production asynchrone, ou client programmatique sans
humain dans la boucle. Enfin, la valeur de l'indexation et du
protocole d'accès aux données n'est pas uniforme mais
\emph{conditionnelle au périmètre} de la demande. Aucune des variantes
comparées n'est universellement préférable, parce qu'aucune ne
s'adresse à la même phase du parcours pédagogique de SprintUp.

Cette pluralité n'est pas une indécision de conception : elle reflète
la diversité réelle des usages que la plateforme doit servir. Elle
déplace toutefois la question d'un niveau. Tant que l'on raisonne
variante par variante, on décrit des formes d'agents isolées ; dès
lors que ces formes coexistent, il faut une logique de niveau
supérieur pour décider laquelle invoquer, pour observer ce qu'elles
font une fois lancées, et pour reconnaître ce qu'elles ne savent pas
encore faire ensemble. Le présent chapitre prend en charge cette
logique. Il ne propose pas \emph{une} orchestration unique qui
remplacerait les variantes ; il propose une \emph{stratégie} qui
organise leur coexistence, en décrit l'observation et en énonce
honnêtement les limites.

Le chapitre s'organise en quatre temps. La première partie définit la
stratégie d'orchestration des variantes : un rappel resserré des trois
surfaces d'interaction héritées du chapitre~2, la politique de routage
qui choisit la variante adaptée à une situation, les sous-graphes que
plusieurs variantes mutualisent, et la composabilité aujourd'hui
limitée de ces surfaces. La deuxième partie décrit ce que signifie
\emph{observer} cette couche d'agents en fonctionnement, à partir de
trois artefacts d'observation et d'un identifiant de fil servant de
trame de corrélation. La troisième partie dresse un bilan lucide des
limites actuelles et des besoins de fiabilisation, au niveau de
l'architecture et non de l'exploitation. La dernière partie fait le
bilan du chapitre et prépare la transition vers le chapitre~4, qui
traitera spécifiquement de la sécurité applicative de cette couche
d'agents.

%% ────────────────────────────────────────────────────────────
\section{Stratégie d'orchestration des variantes}
\label{sec:strategie-orchestration}

La thèse de cette partie tient en une phrase : il n'y a pas une
orchestration de la couche d'agents, mais une stratégie qui choisit la
variante adaptée à la phase du travail et au contexte d'intégration.
Cette stratégie se lit sur trois plans complémentaires. Le premier est
celui des \emph{surfaces} par lesquelles une demande entre dans le
système. Le deuxième est celui du \emph{routage}, c'est-à-dire des
signaux qui déterminent quelle variante répond à une demande donnée.
Le troisième est celui de la \emph{mutualisation}, qui décrit ce que
les variantes partagent en interne pour éviter la duplication. La
partie se termine sur une discussion franche de la composabilité :
ces surfaces fonctionnent aujourd'hui en parallèle, mais ne se
composent pas encore en flux de bout en bout.

\subsection{Trois surfaces d'interaction}
\label{sec:surfaces}

Le chapitre~2 a décrit en détail les variantes de l'agent~3 et les a
relues à travers leurs profils d'intégration (voir la lecture
comparative de la section~\ref{sec:comparatif}). Nous n'en reprenons
ici que la structure d'ensemble, sous l'angle qui importe pour
l'orchestration : chaque variante expose une \emph{surface
d'interaction} distincte, c'est-à-dire une manière propre pour un
acteur — humain ou logiciel — de déclencher une génération et d'en
suivre le déroulement. Trois surfaces se dégagent, et elles ne sont
pas redondantes : chacune répond à une posture humaine que les autres
ne servent pas aussi bien.

\paragraph{Surface conversationnelle.}
La première surface est celle du dialogue. L'enseignant s'adresse à un
agent dans un fil de discussion ; un protocole de streaming lui livre
les artefacts au fur et à mesure de leur production, et les questions
de l'agent prennent la forme d'interruptions structurées. Deux
variantes l'incarnent : la voie conversationnelle à supervision
\gls{llm} avec \emph{Human-in-the-Loop} (la variante
\mcptool{syllabus-generator}, décrite à la section~\ref{sec:v1-op1})
et la voie à sous-agents spécialistes (la variante \mcptool{deepagent}
de la section~\ref{sec:v1-op2}). Toutes deux supposent un enseignant
présent, prêt à clarifier une commande mal spécifiée ou à corriger une
dérive en cours de génération. C'est la surface du
\emph{co-construire à voix haute}.

\paragraph{Surface atelier.}
La deuxième surface renverse cette posture. Plutôt que de dialoguer,
l'enseignant définit d'abord la structure — unités et activités — sous
forme de brouillons, puis déclenche explicitement la génération par
vagues. La variante d'atelier (le couple \mcptool{unit-bulk-supervisor}
et \mcptool{unit-activity-generator}, décrit à la
section~\ref{sec:v2-op1}) ne pratique ni streaming de tokens ni
interruption en ligne : la progression d'une activité n'est exposée
que par une machine d'état publique. C'est la surface de l'\emph{atelier
rédactionnel asynchrone}, où l'opérateur décide quand lancer une vague
et observe ensuite l'avancement à travers des statuts typés.

\paragraph{Surface programmatique.}
La troisième surface retire l'humain de la boucle d'exécution. Un
système client déclenche la génération d'un syllabus par un appel
programmatique, sans interface conversationnelle ni atelier. La
variante d'orchestration par kit de développement (la variante
\mcptool{sdk-orchestrator}, décrite à la section~\ref{sec:v2-op2})
répond par construction à ce contrat : un appel non bloquant lance la
génération en arrière-plan, et la seule observation disponible reste,
là encore, la machine d'état publique. C'est la surface du
\emph{déclenchement industriel}, qui sert une autre application, un
traitement par lots ou une ingestion automatisée.

Il est éclairant de relire ces trois surfaces à travers les axes de
comparaison établis au chapitre~2 (voir la
section~\ref{sec:axes-compar}), car la stratégie d'orchestration en
hérite directement. La \emph{propriété de l'orchestration} les sépare
nettement : sur la surface conversationnelle, un superviseur \gls{llm}
décide à chaque tour du prochain pas, ce qui maximise la souplesse au
prix de la reproductibilité ; sur les surfaces atelier et
programmatique, l'enchaînement des activités est porté par du code
déterministe, ce qui rend l'exécution reproductible mais ferme la porte
à l'adaptation en ligne. La \emph{granularité de l'artefact
inspectable} varie de la même façon : un flux de jetons éphémère sur la
surface conversationnelle, des statuts typés et persistants sur les
deux autres. Quant à la \emph{posture humaine} — enseignant présent,
opérateur asynchrone, client sans humain — elle constitue la différence
la plus visible, et c'est elle que le routage doit reconnaître en
premier. Choisir une surface, c'est donc choisir d'un même geste un
mode d'orchestration, un grain d'observation et une place pour
l'humain ; ces trois choix sont solidaires et ne se recombinent pas
librement.

Ces trois surfaces ne sont pas trois implémentations concurrentes d'un
même besoin : elles répondent à trois moments distincts du travail
pédagogique. La surface conversationnelle sert l'exploration et la
clarification ; la surface atelier sert la production maîtrisée d'un
volume connu ; la surface programmatique sert l'intégration dans un
système tiers. La stratégie d'orchestration ne consiste donc pas à
préférer l'une à l'autre, mais à reconnaître laquelle convient à la
situation présente — ce qui est précisément l'objet du routage.

\subsection{Politique de routage}
\label{sec:routage}

Router, ici, signifie choisir la variante qui répondra à une demande.
Il importe de préciser d'emblée le statut de ce qui suit : nous
décrivons un \emph{cadre de décision}, et non un moteur de routage
automatique. Aujourd'hui, le routage est très majoritairement
\emph{explicite} — c'est l'utilisateur ou le système intégrateur qui
choisit la surface, donc la variante. Le routage assisté, où le
système suggérerait la variante adaptée, relève des perspectives et
n'est pas implémenté. Cette honnêteté de cadrage est essentielle : il
serait trompeur de présenter comme un algorithme sophistiqué ce qui
est, pour l'essentiel, un choix laissé à l'acteur.

Le choix de la variante se laisse néanmoins formaliser comme une
fonction de trois familles de signaux :
\[
  \textsf{routage}(\textit{intention},\ \textit{contexte},\
  \textit{contrat})\ \longrightarrow\ \textit{variante}.
\]
La première famille est l'\textbf{intention déclarée} : ce que
l'acteur veut faire et, indissociablement, la surface qu'il choisit
pour le faire. Un enseignant qui souhaite co-construire un syllabus en
dialoguant ouvre un fil conversationnel ; un enseignant qui veut
produire en bloc une série d'activités déjà structurées passe par
l'atelier. L'intention et la surface sont, dans l'état actuel du
système, presque confondues, et c'est ce qui rend le routage
explicite : déclarer l'intention, c'est déjà choisir la variante.

La deuxième famille est le \textbf{contexte de session}, qui module
un choix à l'intérieur d'une surface donnée. La présence ou l'absence
d'un syllabus déjà lié à la session en est l'exemple le plus net.
Lorsque la génération d'activités s'inscrit dans un syllabus existant,
l'agent outillé décrit au chapitre~2 (l'agent~2, décrit à la
section~\ref{sec:agent2}) est pertinent : il peut consulter les
activités sœurs et ancrer sa production dans le vocabulaire et les
objectifs de l'unité. En l'absence d'un tel rattachement, le contexte
ne fournit rien à ancrer, et l'on retombe naturellement sur le
comportement du générateur sans outillage (l'agent~1, décrit à la
section~\ref{sec:agent1}). Ce basculement reprend directement le
constat de la valeur conditionnelle de l'ancrage établi à la
section~\ref{sec:a1-vs-a2} : le contexte de session détermine s'il y a
lieu de payer le coût de l'ancrage ou non.

La troisième famille est le \textbf{contrat d'intégration}, qui
distingue une sollicitation interactive d'un appel programmatique. Une
interface utilisateur interactive ouvre la voie aux surfaces
conversationnelle et atelier ; un appel programmatique impose, par
construction, la surface programmatique. Le contrat n'est pas un
simple détail d'implémentation : il conditionne la présence ou non
d'un humain susceptible de répondre à une interruption, et donc la
possibilité même d'un dialogue en cours de génération.

\begin{figure}[H]
\centering
\begin{tikzpicture}[
    font=\small,
    sig/.style={rectangle, rounded corners=4pt, draw=myblue!60,
      line width=0.6pt, fill=myblue!6, align=center,
      text width=3.5cm, minimum height=1.1cm, inner sep=4pt},
    dec/.style={diamond, aspect=2.2, draw=myblue!70, line width=0.7pt,
      fill=myblue!10, align=center, text width=2.6cm, inner sep=1pt},
    result/.style={rectangle, rounded corners=4pt, draw=gray!70,
      line width=0.6pt, fill=gray!10, align=center,
      text width=3.4cm, minimum height=1.1cm, inner sep=4pt},
    arrow/.style={->, thick, -{Stealth[length=2.8mm]}, color=black!70}
]
\node[sig] (i) at (0,1.7)  {Intention déclarée\\\scriptsize(et surface choisie)};
\node[sig] (c) at (0,0)    {Contexte de session\\\scriptsize(syllabus lié ?)};
\node[sig] (k) at (0,-1.7) {Contrat d'intégration\\\scriptsize(interactif / programmatique)};
\node[dec] (d) at (5.2,0)  {Cadre de\\décision};
\node[result] (v) at (10.0,0) {Variante retenue\\\scriptsize(surface + agent)};
\draw[arrow] (i.east) -- (d.north west);
\draw[arrow] (c.east) -- (d.west);
\draw[arrow] (k.east) -- (d.south west);
\draw[arrow] (d.east) -- (v.west);
\end{tikzpicture}
\caption{Cadre de décision du routage : trois familles de signaux
  (intention déclarée, contexte de session, contrat d'intégration)
  déterminent la variante invoquée. Le cadre est aujourd'hui appliqué
  de manière explicite par l'acteur ; sa formalisation prépare un
  routage assisté laissé en perspective.}
\label{fig:routage}
\end{figure}

La figure~\ref{fig:routage} résume ce cadre. Il faut en lire la
flèche centrale pour ce qu'elle est : une décision aujourd'hui prise
par un acteur humain ou par un système intégrateur, non par un
composant logiciel dédié. Formaliser le routage comme une fonction
présente malgré tout deux intérêts. D'une part, cela rend explicites
les signaux qui \emph{devraient} gouverner le choix, ce qui aide à
documenter et à enseigner l'usage de la couche d'agents. D'autre part,
cela trace la voie d'un futur routage assisté : un composant pourrait,
à terme, suggérer une variante à partir de l'intention exprimée et du
contexte de session, sans retirer à l'enseignant la décision finale.
Cette évolution est inscrite aux perspectives et n'est pas présentée
ici comme un acquis.

Quelques situations concrètes illustrent l'application de ce cadre. Un
enseignant qui dispose d'une commande encore floue — une discipline, un
niveau, mais aucun plan arrêté — a intérêt à ouvrir un fil
conversationnel, où l'agent l'aidera à clarifier l'intention avant de
produire : l'intention déclarée pointe vers la surface
conversationnelle. Un enseignant qui a déjà saisi une douzaine
d'activités en brouillon et souhaite les voir générées d'un bloc relève
au contraire de la surface atelier, où il déclenchera une vague et en
suivra l'avancement par statuts. Une application tierce qui veut
produire un grand nombre de syllabus sans intervention humaine relève
de la surface programmatique, par contrat d'intégration. Et, à
l'intérieur d'une génération d'activités, le contexte de session
tranche entre l'agent outillé et l'agent sans outillage selon que la
demande s'inscrit ou non dans un syllabus lié — c'est-à-dire selon
qu'il y a, ou non, un existant à consulter et à ne pas dupliquer.

Ces signaux ne sont pas toujours indépendants, et leur articulation
suit une forme de préséance. Le contrat d'intégration prime : un appel
programmatique exclut par construction les surfaces interactives,
quelle que soit l'intention exprimée, car aucun humain n'est disponible
pour répondre à une interruption. Une fois la surface fixée par
l'intention et le contrat, le contexte de session intervient à un grain
plus fin, pour sélectionner l'agent et le degré d'ancrage. Reconnaître
cette préséance n'ajoute aucune automatisation : elle décrit simplement
l'ordre dans lequel les signaux contraignent le choix, et ce que
devrait respecter un futur routage assisté pour rester cohérent avec
l'usage actuel.

\subsection{Sous-graphes communs et mutualisation}
\label{sec:mutualisation}

Si aucune variante n'est universelle, alors le système doit faire
coexister plusieurs formes d'agents. La contrepartie technique de ce
constat est claire : pour éviter de réécrire quatre fois la même
logique, les variantes doivent partager leurs briques internes. La
mutualisation des sous-graphes est le pendant concret de la conclusion
« aucune variante n'est universelle » : la diversité s'exprime au
niveau des surfaces et des postures, pas au niveau de chaque opération
élémentaire.

Trois briques se prêtent particulièrement à ce partage. La première
est le \textbf{sous-graphe de génération d'activité}, qui enchaîne la
planification d'une activité, sa rédaction, puis son indexation dans
la base sémantique. Cette séquence constitue l'unité de travail
élémentaire commune à la surface atelier et à la surface
programmatique : dans les deux cas, le résultat n'est pas un flux
conversationnel mais une activité persistée, dont l'avancement est
exposé par la machine d'état publique. La voie conversationnelle à
sous-agents mobilise la même logique à travers son sous-agent
rédacteur. Extraire ce sous-graphe en une brique réutilisable évite
que chaque surface réimplémente, à sa manière, la même chaîne
planifier-rédiger-indexer, avec le risque de divergences silencieuses
entre variantes.

La deuxième brique est le \textbf{critic pédagogique}, ce sous-agent
dont le rôle est d'examiner une production avant de la considérer comme
acceptable, au regard de la barre de qualité définie pour l'agent
concerné. Mutualiser le critic, c'est garantir qu'une activité produite
par la surface atelier et la même activité produite par la surface
programmatique sont jugées selon les mêmes exigences. C'est aussi
concentrer en un seul point l'évolution de ces exigences : relever la
barre de qualité ou affiner un critère se fait une fois, et bénéficie à
toutes les surfaces qui invoquent le critic.

La troisième brique est l'\textbf{accès aux outils de lecture et
l'indexation}. La consultation des activités existantes, la recherche
par similarité — notamment l'outil \mcptool{find_related_activities} —
et l'écriture indexée constituent un socle partagé, exposé par le
protocole d'accès aux données et soumis au principe du moindre
privilège rappelé à la section~\ref{sec:moindre-privilege}. Ce socle
est ce qui rend l'ancrage possible quel que soit le point d'entrée :
que la demande arrive par dialogue, par atelier ou par appel
programmatique, c'est le même catalogue d'outils, soumis aux mêmes
restrictions, qui sert de pont vers l'existant. Mutualiser cet accès,
c'est aussi mutualiser la discipline de sécurité qui l'encadre — un
point sur lequel le chapitre~4 reviendra.

L'intérêt de cette mutualisation est double. Sur le plan de la
maintenance, factoriser un sous-graphe de rédaction ou un critic
partagé réduit la surface de code à faire évoluer : une amélioration
de l'anti-redondance ou de la barre de qualité profite simultanément
à plusieurs surfaces. Sur le plan de la cohérence, le partage garantit
qu'une même demande, traitée par des surfaces différentes, s'appuie
sur les mêmes garde-fous d'ancrage et de qualité. La diversité des
variantes ne se paie donc pas en divergence de comportement sur les
opérations de base : elle se concentre là où elle apporte de la valeur,
c'est-à-dire dans la posture d'interaction et la propriété de
l'orchestration.

Il faut toutefois mesurer la portée de cette mutualisation. Partager
un sous-graphe n'est pas composer les surfaces entre elles : c'est
réutiliser une brique \emph{à l'intérieur} d'une variante donnée. La
partie sur la composabilité examinera précisément ce que la
mutualisation ne règle pas, à savoir le chaînage des surfaces en un
flux unique. Auparavant, un cas particulier de mutualisation mérite
d'être traité à part, car il touche directement au coût de
fonctionnement : le partage du contexte de session.

\subsection{Mutualisation du contexte et moment de l'ancrage}
\label{sec:cache-contexte}

La lecture comparative du chapitre~2 a laissé une question ouverte,
qu'il revient à ce chapitre de reprendre. L'ancrage d'une activité sur
l'existant a un coût observable — de l'ordre de plusieurs dizaines de
pour cent de latence supplémentaire et un volume de jetons d'entrée
sensiblement accru — et ce coût est entièrement imputable à la phase
de lecture qui précède la génération (voir la
section~\ref{sec:a1-vs-a2}). La question n'est pas de savoir s'il faut
payer ce coût, puisque l'ancrage conditionne la non-redondance ; elle
est de savoir \emph{quand} le payer : à chaque génération, ou une fois
pour un ensemble de générations relevant d'une même session.

Dans l'état actuel de la conception, la lecture d'ancrage est
recalculée \emph{à chaque génération}. Ce choix est simple et correct :
chaque activité produite consulte l'existant au moment où elle est
produite, ce qui garantit qu'elle s'appuie sur l'état le plus récent du
syllabus. Il a néanmoins une conséquence visible lorsqu'une même
session enchaîne plusieurs générations sur le même syllabus — cas
fréquent sur la surface atelier : le contexte de lecture, pour
l'essentiel identique d'une activité à l'autre, est reconstitué autant
de fois qu'il y a de générations, et le tarif de l'ancrage est acquitté
à répétition.

Le levier architectural correspondant consiste à \emph{mutualiser le
contexte de session} : lire une fois le contexte pertinent d'un
syllabus — activités sœurs, plan d'ensemble, objectifs — puis le
réutiliser pour les générations successives d'une même session, plutôt
que de le reconstruire à chaque fois. Ce levier prolonge directement la
logique de routage exposée plus haut : on ne paie l'ancrage que
lorsqu'il apporte de la valeur, et, idéalement, on ne le paie qu'une
fois par session cohérente plutôt qu'une fois par activité. Il s'agit
d'une mise en cache du contexte au niveau de la couche d'agents, et non
d'un mécanisme de stockage particulier : la notion reste ici
architecturale.

Un tel partage soulève en retour une question de fraîcheur. Un contexte
mutualisé peut se périmer si des activités sont créées ou modifiées en
cours de session ; sa réutilisation doit donc être bornée par une
invalidation, naturellement adossée aux événements qui font évoluer le
syllabus — transitions de la machine d'état des activités, nouvelles
indexations. L'identifiant de fil introduit pour l'observabilité offre
d'ailleurs une portée naturelle à ce contexte partagé : il délimite la
session à l'intérieur de laquelle le contexte peut être réutilisé sans
risque d'incohérence. Nous présentons ce levier comme une
\emph{perspective} cohérente avec la stratégie d'orchestration, et non
comme un mécanisme déjà en place : la conception actuelle recalcule
l'ancrage, et la mutualisation du contexte de session reste à
concevoir.

\subsection{Composabilité limitée}
\label{sec:composabilite}

Une lecture honnête de la conception actuelle impose de reconnaître
une limite : les surfaces sont \emph{parallèles mais non composées}.
Chacune offre un chemin complet, de la demande au résultat, mais
aucune ne s'enchaîne automatiquement à une autre. Deux exemples
concrets le montrent. Un déclenchement par la surface programmatique
produit un syllabus complet, mais n'ouvre pas ensuite, de lui-même, un
atelier de révision interactive : si l'enseignant souhaite reprendre
le résultat à la main, il doit le rouvrir explicitement dans la
surface atelier. Symétriquement, une session conversationnelle ne
verse pas automatiquement son résultat dans un atelier de production
par lots : le passage d'une surface à l'autre reste un geste manuel.

La cause de cette limite est structurelle plus qu'accidentelle :
il n'existe pas, à ce jour, de graphe global unique qui porterait
l'état d'un travail à travers les surfaces. Chaque surface gère son
propre état — le fil de conversation pour les voies interactives, la
machine d'état des activités pour les voies atelier et programmatique —
sans qu'un orchestrateur de niveau supérieur ne relie ces états en une
trajectoire continue. La conséquence est qu'un travail entamé sur une
surface ne peut être prolongé sur une autre que par une intervention
explicite, qui recommence depuis l'état observable de la surface
d'arrivée.

Nous présentons ce point comme une limite assumée, et non comme une
promesse non tenue. Composer les surfaces en flux de bout en bout —
par exemple, déclencher programmatiquement une génération puis basculer
en atelier de révision sans perte d'état — supposerait un orchestrateur
global et un état partagé que la conception actuelle ne fournit pas. Ce
besoin est réel et il est inscrit aux perspectives ; il rejoint la
discussion des limites de la partie consacrée à la fiabilisation, dans
la troisième partie de ce chapitre. Le reconnaître clairement vaut
mieux que de laisser croire à une composabilité qui n'existe pas
encore.

On peut esquisser, sans entrer dans l'implémentation, ce qu'exigerait
une telle composition. Il faudrait d'abord un \emph{état partagé} qui
survive au passage d'une surface à l'autre, de sorte qu'un travail
entamé programmatiquement conserve son identité et son contenu
lorsqu'il est repris en atelier. Il faudrait ensuite un \emph{contrat
de transfert} explicite entre surfaces, décrivant ce qu'une surface
remet à la suivante et dans quel état le travail doit se trouver pour
être repris sans ambiguïté. Il faudrait enfin un \emph{orchestrateur de
niveau supérieur} capable d'enchaîner ces transferts en respectant les
postures humaines de chaque surface — par exemple, ne basculer en
révision interactive que si un humain est effectivement disponible. Ces
trois éléments dessinent un programme de travail cohérent ; la
conception actuelle, qui privilégie des surfaces autonomes et bien
délimitées, en constitue une fondation raisonnable plutôt qu'un
obstacle.

\subsection{Synthèse : une stratégie, non une orchestration unique}
\label{sec:synthese-orchestration}

Les quatre plans qui précèdent dessinent ensemble une posture de
conception cohérente, qu'il convient de résumer avant d'aborder
l'observabilité. La stratégie d'orchestration ne cherche pas à fondre
les variantes en un orchestrateur universel : elle les fait coexister
en reconnaissant que chacune sert une posture d'usage distincte. Les
\emph{surfaces} encapsulent ces postures ; le \emph{routage}, explicite
aujourd'hui, choisit la surface et la variante à partir de l'intention,
du contexte et du contrat ; la \emph{mutualisation} garantit que cette
diversité de surfaces ne se paie pas en duplication des opérations de
base ni en divergence des garde-fous de qualité et d'ancrage ; la
\emph{composabilité}, enfin, demeure limitée et constitue le principal
chantier ouvert.

Cette posture a une vertu méthodologique : elle distingue nettement ce
qui relève d'un choix de conception assumé de ce qui relève d'une
limite à lever. Faire coexister plusieurs surfaces est un choix, motivé
par la pluralité réelle des usages de SprintUp. Ne pas les composer
encore est une limite, dont la résolution suppose un état partagé et un
orchestrateur de niveau supérieur. Tenir ces deux registres séparés —
le choix et la limite — est précisément ce qui permet de défendre la
conception sans la survendre. C'est aussi ce qui prépare la partie
suivante : une stratégie que l'on assume doit pouvoir être observée, de
sorte que ses choix comme ses limites soient vérifiables en
fonctionnement, et non seulement affirmés sur le papier.

%% ────────────────────────────────────────────────────────────
\section{Observabilité de la couche d'agents}
\label{sec:observabilite}

Une stratégie d'orchestration n'est défendable que si l'on peut
observer ce qu'elle produit en fonctionnement. Observer, ici, ne
signifie pas déployer un outillage de supervision particulier — le
choix d'un tel outillage sort du périmètre de la couche d'agents — mais
identifier \emph{ce qui} est observable, à quel niveau d'abstraction,
et comment relier ces observations entre elles. Cette partie décrit
trois artefacts d'observation indépendants de toute technologie,
montre comment un identifiant de fil les relie, énumère les mesures
opérationnelles utiles à suivre dans le temps, et propose un schéma de
diagnostic par remontée. Conformément à la discipline adoptée dans
tout le chapitre, nous distinguons à chaque étape ce qui est
\emph{déjà conçu et en place} de ce qui reste \emph{à raffiner}.

\subsection{Trois artefacts d'observation}
\label{sec:artefacts}

Trois artefacts d'observation coexistent et se complètent. Pris
séparément, chacun ne donne qu'une vue partielle ; c'est leur
corrélation qui rend la couche d'agents intelligible.

\paragraph{Traces d'exécution.}
Le premier artefact est la trace d'exécution émise par les
superviseurs et les sous-agents. Une telle trace est une séquence de
phases et d'appels d'outils : sollicitation d'un outil de lecture,
délégation à un sous-agent, passage du planificateur au rédacteur,
appel du critic, et ainsi de suite. Ce niveau de détail est
\emph{en place} : la conception des variantes, en particulier les
voies à supervision \gls{llm}, produit naturellement ces traces, et
elles sont la seule fenêtre sur le raisonnement interne des surfaces
conversationnelles, qui n'exposent pas de machine d'état publique. La
trace est l'artefact le plus riche, mais aussi le plus volumineux et
le moins structuré : elle se lit a posteriori, pour comprendre
\emph{comment} un résultat a été obtenu.

La richesse de la trace est à la fois sa force et sa difficulté. Elle
permet de reconstituer le cheminement exact d'une variante — quel
outil a été appelé, dans quel ordre, avec quel effet — ce qui est
indispensable pour les surfaces conversationnelles, où nul état public
ne résume l'avancement. Mais cette richesse a un coût d'interprétation :
une trace n'est pas un résumé, c'est un journal détaillé, et en tirer
un diagnostic suppose de savoir quoi y chercher. C'est pourquoi la
trace gagne à être lue \emph{en regard} des deux autres artefacts
plutôt qu'isolément, comme la suite de cette partie le montre.

\paragraph{Machine d'état publique des activités.}
Le deuxième artefact est la machine d'état publique des activités de
la Version~2, déjà décrite et illustrée au chapitre~2 (voir la
figure~\ref{fig:v2-status}). Une activité y progresse à travers une
suite d'états — \mcptool{draft}, \mcptool{queued}, \mcptool{searching},
\mcptool{generating}, \mcptool{ready} — assortie d'états terminaux
d'échec — \mcptool{failed} et \mcptool{blocked}. Cet artefact présente
une propriété décisive pour l'observabilité : il fait partie du
\emph{contrat public} du système. Un client peut suivre l'avancement
d'une génération en lisant uniquement ces états, sans interpréter le
moindre protocole de streaming ni accéder au raisonnement interne de
l'agent. La machine d'état est ainsi l'artefact le plus pauvre en
détail mais le plus stable et le plus universellement consommable ;
elle est \emph{en place} pour les voies atelier et programmatique.

Cette pauvreté apparente est en réalité une vertu d'interface. Parce
qu'elle ne dépend ni du modèle de langage employé, ni du détail de
l'orchestration interne, la machine d'état offre une surface
d'observation stable dans le temps : un client construit sur ces états
ne se brise pas lorsque le raisonnement interne d'une variante évolue.
Elle constitue ainsi le point d'observation de référence pour quiconque
consomme la couche d'agents sans en connaître les rouages — et, à ce
titre, le premier artefact que l'on consulte lors d'un diagnostic.

\paragraph{Mesures opérationnelles.}
Le troisième artefact est l'ensemble des mesures opérationnelles —
latence, volume de jetons, fréquence d'erreurs — qui caractérisent une
exécution indépendamment de la variante. Contrairement aux deux
précédents, cet artefact relève aujourd'hui surtout du
\emph{préconisé} : l'architecture permet de produire et de collecter
ces grandeurs, mais leur agrégation systématique dans un dispositif de
suivi continu reste à construire. Nous le présentons donc comme un
\emph{dispositif d'observation à mettre en place}, et non comme un
système de supervision déjà déployé. Les mesures elles-mêmes sont
détaillées plus loin, dans la partie consacrée aux grandeurs utiles.

La complémentarité de ces trois artefacts tient à leurs niveaux
d'abstraction. La machine d'état répond à la question « où en est ce
travail ? » ; la trace répond à « par quel cheminement en est-il
arrivé là ? » ; les mesures répondent à « à quel coût, et dans quelles
marges de temps ? ». Aucun ne se substitue aux autres, et c'est leur
mise en regard qui permet de raisonner sur l'architecture.

\subsection{Corrélation par identifiant de fil}
\label{sec:correlation}

Pour relier ces trois artefacts, encore faut-il une clé commune. Cette
clé est un \emph{identifiant de fil} (ou identifiant de conversation),
attribué à chaque travail entrant dans le système. Le même
identifiant accompagne la demande de l'acteur, la trace d'exécution
émise par les agents, et les transitions d'état des activités
produites. Il joue le rôle de \emph{fil conducteur} : à partir d'un
seul identifiant, on peut remonter d'un symptôme observable jusqu'au
détail du raisonnement qui l'a produit.

\begin{figure}[H]
\centering
\begin{tikzpicture}[
    font=\small,
    lane/.style={font=\small\bfseries, align=right, text width=3.0cm},
    ev/.style={rectangle, rounded corners=3pt, draw=myblue!60,
      line width=0.5pt, fill=myblue!8, align=center,
      font=\footnotesize, minimum height=0.8cm, inner sep=3pt},
    st/.style={rectangle, rounded corners=3pt, draw=gray!70,
      line width=0.5pt, fill=gray!12, align=center,
      font=\footnotesize\ttfamily, minimum width=1.5cm,
      minimum height=0.7cm, inner sep=2pt},
    axis/.style={->, thick, -{Stealth[length=2.6mm]}, color=black!55},
    link/.style={dashed, color=black!35, line width=0.4pt}
]
%% --- trois couloirs ---
\node[lane] (l1) at (-0.3,2.4) {Demande\\acteur};
\node[lane] (l2) at (-0.3,1.0) {Trace\\d'exécution};
\node[lane] (l3) at (-0.3,-0.6) {Machine\\d'état};

%% --- couloir 1 : demande ---
\node[ev] (req) at (2.2,2.4) {requête};

%% --- couloir 2 : trace LLM ---
\node[ev] (t1) at (2.2,1.0) {intake};
\node[ev] (t2) at (4.6,1.0) {recherche};
\node[ev] (t3) at (7.0,1.0) {rédaction};
\node[ev] (t4) at (9.4,1.0) {critic};

%% --- couloir 3 : états ---
\node[st] (s1) at (2.2,-0.6) {queued};
\node[st] (s2) at (4.6,-0.6) {searching};
\node[st] (s3) at (7.0,-0.6) {generating};
\node[st] (s4) at (9.7,-0.6) {ready};

%% --- axe du temps ---
\draw[axis] (1.0,-1.6) -- (10.6,-1.6) node[right, font=\footnotesize] {temps};

%% --- liens verticaux de corrélation (même thread) ---
\draw[link] (t1) -- (s1);
\draw[link] (t2) -- (s2);
\draw[link] (t3) -- (s3);
\draw[link] (t4.south) -- (s4.north);
\draw[link] (req) -- (t1);

\node[font=\footnotesize\itshape, color=black!60, align=center]
  at (5.8,3.2) {un même identifiant de fil relie les trois couloirs};
\end{tikzpicture}
\caption{Chronologie d'un fil. Une demande unique se déploie sur trois
  plans corrélés par un même identifiant de fil : la requête de
  l'acteur, la trace d'exécution (phases et appels d'outils) et la
  machine d'état publique des activités. La superposition permet de
  rattacher chaque état observable à la phase qui l'a produit.}
\label{fig:thread-timeline}
\end{figure}

La figure~\ref{fig:thread-timeline} illustre cette corrélation sous la
forme d'une chronologie. Trois couloirs y avancent en parallèle le
long d'un même axe de temps. Le premier porte la demande de l'acteur.
Le deuxième déroule la trace d'exécution — intake, recherche,
rédaction, critic — c'est-à-dire le raisonnement interne de la
variante. Le troisième suit la machine d'état publique, dont les
transitions (\mcptool{queued}, \mcptool{searching}, \mcptool{generating},
\mcptool{ready}) reflètent, vues de l'extérieur, l'avancement du
travail. L'identifiant de fil, partagé par les trois couloirs, autorise
la lecture verticale : à un état public donné correspond une phase de
la trace, et à cette phase un coût et une durée mesurables. C'est cette
lecture verticale qui transforme trois flux d'information séparés en un
récit unique et exploitable d'un travail.

Il convient de situer le statut de ce dispositif. L'existence de
traces détaillées au niveau des agents et celle d'une machine d'état
publique sont \emph{en place} : ce sont des propriétés de conception
des variantes décrites au chapitre~2. La généralisation d'un
identifiant de fil unique comme clé de corrélation systématique, à
travers tous les artefacts et toutes les surfaces, relève en revanche
du \emph{préconisé} : le principe est posé, mais son application
uniforme demande un travail d'instrumentation qui reste à mener.

\subsection{Mesures opérationnelles utiles}
\label{sec:metriques}

Au-delà des artefacts qualitatifs, quelques grandeurs méritent d'être
suivies dans le temps, car elles éclairent les compromis propres à
chaque variante. L'argument n'est pas dans leur valeur absolue : il est
dans leur \emph{évolution} et dans ce qu'elle révèle des arbitrages de
conception. Nous en retenons quatre.

\begin{itemize}
  \item \textbf{Latence de bout en bout, par variante.} Le temps qui
    sépare le déclenchement d'un travail de la disponibilité de son
    résultat. Cette grandeur ne se compare pas d'une variante à
    l'autre dans l'absolu : une surface conversationnelle vise une
    attente de l'ordre de la minute et tolère le streaming, là où une
    surface programmatique répond de manière non bloquante et se
    qualifie plutôt par le débit d'une vague. La latence éclaire
    surtout l'adéquation d'une variante à l'enveloppe d'usage qu'elle
    sert, telle qu'elle a été décrite à la section~\ref{sec:observations}.
  \item \textbf{Volume de jetons consommés par exécution.} La taille
    du contexte mobilisé pour produire un résultat. Le chapitre~2 a
    montré, à titre d'ordre de grandeur, que l'ancrage d'une activité
    sur l'existant accroît sensiblement le volume de jetons d'entrée
    par rapport à une génération non ancrée (voir la
    section~\ref{sec:a1-vs-a2}). Suivre cette grandeur, c'est mesurer
    le \emph{tarif de l'ancrage} et vérifier qu'il n'est payé que
    lorsqu'il apporte de la valeur.
  \item \textbf{Fréquence et nature des erreurs d'accès aux outils.}
    Les délais dépassés et les échecs d'appel aux outils de lecture ou
    d'écriture. Cette grandeur renseigne sur la robustesse de la
    couche d'accès aux données et, indirectement, sur la part des
    échecs de génération imputable non au raisonnement de l'agent mais
    à l'indisponibilité d'une ressource.
  \item \textbf{Temps passé par une activité dans chaque état.} La
    durée de séjour d'une activité dans les états de la machine
    publique — combien de temps en \mcptool{queued}, en
    \mcptool{searching}, en \mcptool{generating}. Cette grandeur
    localise les goulets d'étranglement le long du flux nominal et
    distingue, par exemple, une attente d'ordonnancement d'une phase
    de recherche anormalement longue.
\end{itemize}

Ces grandeurs reprennent et prolongent les ordres de grandeur de la
lecture comparative du chapitre~2. Il faut en rappeler le statut :
l'objectif n'est pas de produire un classement chiffré des variantes,
mais de \emph{comprendre les compromis} — latence contre ancrage,
débit contre contrôlabilité, richesse de la trace contre stabilité du
contrat public. Suivre ces mesures dans le temps est précisément ce
qui permettrait de faire de ces compromis un objet de décision, et non
une simple intuition de conception.

Deux précautions encadrent l'usage de ces grandeurs. La première est
de ne pas confondre une mesure ponctuelle avec une tendance : une
latence observée sur quelques exécutions ne dit rien de fiable tant
qu'elle n'est pas suivie dans la durée, et c'est précisément le suivi
continu — aujourd'hui à mettre en place — qui donnerait son sens à ces
chiffres. La seconde est de relier chaque grandeur à l'artefact qui
l'explique : une latence anormale renvoie à la trace pour en localiser
la phase coûteuse, un temps de séjour excessif dans un état renvoie à
la machine d'état pour en circonscrire l'étape. Les mesures, en somme,
ne se suffisent pas à elles-mêmes : elles signalent où regarder, et ce
sont les deux autres artefacts qui disent quoi corriger.

\subsection{Diagnostic par remontée}
\label{sec:diagnostic}

Les trois artefacts et leur clé de corrélation prennent tout leur sens
dans une situation concrète : le diagnostic d'un symptôme. Le schéma
proposé est une remontée à trois étages, de la surface observable vers
la cause architecturale. Prenons un symptôme typique : une activité
demeure longtemps dans l'état \mcptool{queued}.

Le premier étage consiste à \textbf{inspecter l'état public}. La
machine d'état indique sans ambiguïté que l'activité n'a pas quitté la
file d'attente : le travail n'a pas échoué, il n'a pas non plus
progressé. Cette observation, à elle seule, oriente déjà le diagnostic
en écartant un échec franc (\mcptool{failed}) ou un blocage propagé
par dépendance (\mcptool{blocked}).

Le deuxième étage consiste à \textbf{inspecter la trace du fil}
correspondant. En suivant l'identifiant de fil, on relie l'activité
bloquée à la séquence de phases qui la concerne, et l'on observe à
quel moment la progression s'interrompt : la phase de recherche
a-t-elle été lancée puis suspendue ? un appel d'outil est-il resté
sans réponse ? l'ordonnancement a-t-il simplement différé le passage à
l'état suivant ?

Le troisième étage consiste à \textbf{raisonner sur l'architecture}.
Selon ce que révèle la trace, l'explication se rattache à un choix de
conception identifiable : un sous-graphe de recherche
particulièrement long sur un syllabus volumineux, un plafond de
parallélisme imposé par l'ordonnanceur qui retient les activités en
file, ou un délai d'accès à un outil de lecture. Le diagnostic ne
s'arrête pas à constater le symptôme : il le rattache à une propriété
de l'architecture, ce qui ouvre la voie à une correction raisonnée.

Le même patron s'applique à d'autres symptômes. Supposons une hausse
de la fréquence des erreurs d'accès aux outils sur une période donnée.
L'état public, là encore, fournit le premier indice : les activités
concernées basculent en échec à la phase de recherche plutôt que de
progresser. La trace du fil précise ensuite la nature de l'erreur — un
délai dépassé, un appel resté sans réponse, plutôt qu'un refus de
raisonnement de l'agent. Le raisonnement architectural conclut enfin
en distinguant ce qui relève de la couche d'accès aux données de ce qui
relève de la logique de l'agent : un échec concentré sur un même outil
de lecture oriente vers la robustesse de cet accès, et non vers une
révision du prompt de l'agent. La valeur du patron tient à cette
capacité de \emph{séparer les causes} — raisonnement, accès aux
données, ordonnancement — au lieu de les confondre dans un symptôme
global.

Ce schéma reste, à ce stade, un \emph{patron conceptuel} : il décrit
la manière de raisonner, non un tableau de bord prêt à l'emploi. La
capacité d'inspecter l'état public et la trace d'un fil est en place ;
l'outillage qui rendrait cette remontée immédiate et systématique —
recherche par identifiant, mise en regard automatique des trois
artefacts — reste à construire. C'est l'une des limites que la partie
suivante examine.

\subsection{Ce que l'observabilité permet de défendre}
\label{sec:synthese-observabilite}

Il vaut la peine de récapituler ce que ce modèle d'observabilité
autorise, en maintenant la distinction entre l'acquis et le projeté.
Sont \emph{en place}, parce qu'ils découlent directement de la
conception décrite au chapitre~2 : l'existence de traces détaillées au
niveau des superviseurs et des sous-agents ; la machine d'état publique
des activités comme contrat d'observation stable ; et la possibilité,
pour un fil donné, de mettre en regard un état public et la phase de
trace correspondante. Ces propriétés suffisent à raisonner, au cas par
cas, sur un comportement observé et à en remonter la cause.

Relèvent en revanche du \emph{préconisé} : la généralisation d'un
identifiant de fil unique comme clé de corrélation systématique ; la
collecte continue et l'agrégation des mesures opérationnelles ; et
l'outillage qui transformerait le diagnostic par remontée en un geste
immédiat plutôt qu'en une investigation manuelle. La frontière entre
ces deux registres n'est pas une faiblesse à dissimuler : c'est la
description exacte de ce sur quoi une soutenance peut s'appuyer sans
exagération. Affirmer que la couche d'agents \emph{peut} être observée,
et préciser à quel prix elle le serait de manière systématique, est
plus solide que de revendiquer une supervision qui n'existe pas. Cette
honnêteté de cadrage conduit naturellement à l'examen des limites, qui
fait l'objet de la partie suivante.

%% ────────────────────────────────────────────────────────────
\section{Limites actuelles et besoins de fiabilisation}
\label{sec:limites}

Les deux premières parties de ce chapitre ont décrit une stratégie
d'orchestration et un dispositif d'observabilité. Il serait malhonnête
de les présenter comme un système pleinement mûr. Cette partie dresse,
au niveau de l'architecture et non de l'exploitation, un bilan lucide
des limites actuelles et des besoins de fiabilisation. Trois familles
de limites se dégagent : celles du routage et de la composition,
celles de la gestion d'état et de la reprise, et celles de
l'observabilité et des métriques.

\subsection{Limites du routage et de la composition}
\label{sec:limites-routage}

La première limite a déjà été énoncée dans la partie sur
l'orchestration, et il convient de la reprendre comme telle. Le
routage demeure \emph{explicite et manuel} : aucun composant ne
combine automatiquement les signaux d'intention, de contexte et de
contrat pour suggérer ou choisir une variante. Le cadre de décision
formalisé à la section~\ref{sec:routage} décrit comment le choix
\emph{devrait} se faire, mais c'est aujourd'hui l'acteur — humain ou
système intégrateur — qui l'exécute. Un routage assisté, qui
proposerait une variante sans confisquer la décision finale, reste une
perspective.

La seconde limite est celle de la composition. Les surfaces —
conversationnelle, atelier, programmatique — ne se chaînent pas en
flux à plusieurs étapes. Il n'existe pas d'orchestrateur global qui
enchaînerait, de manière sûre, un déclenchement programmatique suivi
d'une révision en atelier, ou une exploration conversationnelle suivie
d'une production par lots. Comme nous l'avons noté à la
section~\ref{sec:composabilite}, la cause est l'absence d'un graphe
unique portant l'état d'un travail à travers les surfaces. Tant que
cet état reste cloisonné par surface, la composition demeure un geste
manuel plutôt qu'une trajectoire orchestrée.

Ces deux limites se rejoignent dans un même manque : celui d'un niveau
d'orchestration qui se situerait \emph{au-dessus} des variantes plutôt
qu'à l'intérieur de chacune. Un routage assisté et une composition de
surfaces sont en effet deux facettes d'un même besoin — décider, à
partir de signaux, quelle variante invoquer, puis enchaîner plusieurs
invocations en une trajectoire cohérente. Les énoncer ensemble n'est
pas un aveu de faiblesse : c'est délimiter avec précision le périmètre
de ce qui a été conçu (des surfaces autonomes et solides) et de ce qui
reste à concevoir (leur coordination de niveau supérieur).

\subsection{Limites de la gestion d'état et de la reprise}
\label{sec:limites-etat}

La fiabilité d'une couche d'agents se mesure aussi à sa capacité de
reprendre un travail interrompu. Sur ce plan, la conception actuelle
offre des possibilités réelles mais inégales selon la surface, qu'il
faut décrire sans les surestimer.

Pour les voies conversationnelles, l'état d'un travail est porté par
l'état du graphe de dialogue : une session interrompue peut, en
principe, être reprise à partir de l'état conservé du fil, puisque le
raisonnement progresse par tours successifs dont l'état est
explicitement maintenu. Pour les voies atelier et programmatique, la
reprise s'appuie sur la machine d'état des activités : une activité en
échec peut être remise en file pour une nouvelle tentative, et la
propagation d'un échec aux activités dépendantes est matérialisée par
un état terminal dédié, comme l'illustre la figure~\ref{fig:v2-status}.
La granularité de la reprise est ici l'activité : on relance ce qui a
échoué, sans nécessairement rejouer ce qui avait réussi.

Ces mécanismes ne constituent pas pour autant une reprise
automatique et complète. En particulier, il n'existe pas de mécanisme
de \emph{retour arrière} automatique pour une chaîne d'actions
complexe : annuler proprement un ensemble d'effets partiellement
appliqués, lorsqu'un travail composite échoue à mi-parcours, n'est pas
pris en charge de bout en bout et resterait à concevoir. La reprise
actuelle est, pour l'essentiel, une reprise \emph{en avant} — relancer
ce qui a échoué — et non une annulation transactionnelle de ce qui a
déjà été produit. Ce constat n'est pas un défaut surprenant pour un
travail de cette nature ; il délimite simplement ce qui relève de
l'ingénierie ultérieure.

\subsection{Limites de l'observabilité et des métriques}
\label{sec:limites-observabilite}

La dernière famille de limites concerne l'observabilité elle-même. Il
faut le reconnaître nettement : l'observabilité actuelle n'est pas
\emph{industrialisée}. Il n'existe pas de système complet et automatisé
de supervision, avec tableaux de bord agrégés et alertes, qui
suivrait en continu les grandeurs énumérées à la
section~\ref{sec:metriques}. L'architecture rend ce suivi
\emph{possible} — les artefacts existent, la clé de corrélation est
identifiée — mais le travail qui le rendrait systématique reste à
faire.

Le mot « industrialiser » mérite d'être précisé, pour ne pas déborder
sur un terrain qui n'est pas celui de ce travail. Il ne s'agit pas ici
de choisir un outillage de supervision ni d'en décrire l'exploitation :
ces questions relèvent de l'ingénierie de mise en {\oe}uvre et non de
la conception de la couche d'agents. Il s'agit, plus modestement, de
constater que les artefacts d'observation existent à l'état de
\emph{capacité} — on peut, au cas par cas, lire un état, suivre une
trace, calculer une mesure — sans qu'une collecte régulière et une mise
en regard automatique n'en fassent encore un instrument de pilotage.
Combler cet écart est un travail d'instrumentation identifié, pas une
lacune de conception.

Cette limite n'est pas isolée : elle prolonge celle du diagnostic par
remontée, dont nous avons dit qu'il demeure un patron conceptuel
faute d'outillage dédié. Elle prépare aussi une articulation naturelle
avec le chapitre~4. La journalisation, l'auditabilité et la traçabilité
ne sont pas seulement des instruments d'exploitation : ce sont aussi
des exigences de sécurité. Savoir qui a déclenché quelle variante, sur
quel périmètre, et pouvoir le retracer de manière fiable, relève autant
de l'observabilité que du contrôle d'accès. Industrialiser
l'observabilité et instruire la sécurité applicative sont ainsi deux
chantiers qui se renforcent mutuellement, et le second fera l'objet du
chapitre suivant.

En somme, ce que ce chapitre décrit est une \emph{base architecturale
solide} : des surfaces clairement caractérisées, un cadre de décision
explicite, des artefacts d'observation complémentaires et une clé de
corrélation identifiée. Ce n'est pas un système entièrement mûri et
durci pour une exploitation à grande échelle — ce qui est normal, et
même attendu, pour un projet de fin d'études. Nommer cette distinction
n'affaiblit pas la conception : c'est au contraire la condition d'une
défense honnête de ses apports comme de ses limites.

%% ────────────────────────────────────────────────────────────
\section{Bilan du chapitre et transition vers la sécurité applicative}
\label{sec:ch3-bilan}

Ce chapitre a prolongé la lecture comparative du chapitre~2 sur trois
plans. Là où le chapitre~2 comparait des formes d'agents prises une à
une, ce chapitre en a organisé la coexistence.

Le premier apport est une \textbf{vue stratégique} de la manière dont
les variantes coexistent et sont choisies. En distinguant trois
surfaces d'interaction, en formalisant le routage comme une fonction
de l'intention, du contexte et du contrat, et en identifiant les
sous-graphes mutualisés, le chapitre a montré que la diversité des
variantes n'est pas un éparpillement mais une réponse organisée à des
postures d'usage distinctes. Il a aussi reconnu, sans détour, la
composabilité encore limitée de ces surfaces.

Le deuxième apport est un \textbf{modèle conceptuel d'observabilité},
fondé sur trois artefacts — traces d'exécution, machine d'état publique,
mesures opérationnelles — reliés par un identifiant de fil. Ce modèle
fournit un langage pour raisonner sur le comportement de la couche
d'agents en fonctionnement, depuis le suivi d'un travail jusqu'au
diagnostic d'un symptôme par remontée.

Le troisième apport est une \textbf{vue honnête des limites} et des
besoins de fiabilisation : routage explicite et non assisté, surfaces
parallèles mais non composées, reprise en avant plutôt
qu'annulation transactionnelle, observabilité possible mais non encore
industrialisée. Ces limites tracent, en creux, le programme de travail
d'une version ultérieure.

Au-delà de ces trois apports, le chapitre a maintenu une discipline
transversale qui constitue en elle-même une contribution
méthodologique : pour chaque surface, chaque artefact d'observation et
chaque mécanisme de reprise, il a distingué ce qui est \emph{en place}
de ce qui est \emph{préconisé}. Cette séparation n'est pas un simple
souci de prudence rédactionnelle. Elle reflète la posture attendue d'un
travail d'ingénierie : ne revendiquer comme acquis que ce qui est
effectivement conçu et défendable, et formuler le reste comme une
perspective argumentée. Appliquée à une couche d'agents — où il est
facile de confondre ce que la technologie \emph{permettrait} avec ce
qui est \emph{réalisé} — cette discipline est aussi une garantie
d'honnêteté vis-à-vis du lecteur comme du jury.

Jusqu'ici, le rapport a suivi une progression continue. Les
chapitres~1 et~2 ont posé la conception logique de la couche d'agents :
fondations théoriques de l'orchestration et du protocole d'accès aux
données, puis catalogue comparé des formes d'agents. Le présent
chapitre y a ajouté la stratégie d'orchestration et le dispositif
d'observabilité. Un axe demeure cependant à approfondir, que ce
chapitre n'a abordé qu'en filigrane : la \textbf{sécurité applicative}
de cette couche d'agents.

Cet axe mérite un traitement à part entière. Le principe du moindre
privilège rappelé au chapitre~2 (voir la
section~\ref{sec:moindre-privilege}) constitue un point de départ, mais
il ne dit rien de \emph{qui} sollicite un agent ni de la manière dont
ce contrôle se combine à une gestion des accès par rôle. De même, une
couche d'agents fondée sur de grands modèles de langage présente des
surfaces d'attaque spécifiques — détournement par instructions
injectées, fuite de données par les réponses — dont la mitigation
n'a été ici qu'effleurée comme préoccupation, sans analyse dédiée. Le
chapitre~4 prendra en charge cet axe. Il examinera comment le moindre
privilège du protocole d'accès s'articule à un contrôle d'accès fondé
sur les rôles, comment les injections d'instructions et les fuites de
données se mitigent dans ce type d'architecture, et comment raisonner
sur les contraintes de conformité dans ce contexte. Le chapitre
suivant approfondit cet axe en examinant, sous l'angle de la
cybersécurité, les surfaces d'attaque propres à une telle couche
d'agents et les mesures de mitigation envisageables dans le cadre de
SprintUp.

\newpage
